\section{Translate characters into a computation}\label{expression-purpose}
An expression gives the evaluator two kinds of information. Digit sequences supply
values; operators and parentheses specify how to combine those values. Reading the
characters in order does not imply applying every operator immediately. In
\texttt{18 - 3 * (2 + -4)}, the subtraction's right operand includes an entire
multiplication. That multiplication must obtain the parenthesized value before the
subtraction can finish. The final result is 24.

The parser will recognize the expression and calculate its value in the same pass.
It will not allocate a syntax tree. A returned value represents the expression that
one parser call has consumed, while a shared cursor identifies the next unconsumed
byte. The caller of that parser routine decides whether that next byte belongs to
another operator at its own level, closes a group, or ends the complete input.

This division of work makes a recursive call useful. A group can ask the expression
routine to evaluate its contents even while an outer expression is unfinished.
The unfinished operation must survive that call. In assembly, that means preserving
both the return address and any live local values that the nested call can overwrite.
The stack will hold those separate invocations; the input itself remains read-only.
